\documentclass[12pt,a4paper]{article}
\usepackage[margin=25mm, headheight=15pt, headsep=10pt]{geometry}
\usepackage{fontspec}
\usepackage[table]{xcolor} % Note: table option is required for rowcolors
\usepackage{titlesec}
\usepackage{tocloft}
\usepackage{fancyhdr}
\usepackage{booktabs}
\usepackage{tcolorbox}
\tcbuselibrary{skins, breakable}
\usepackage[colorlinks=true, linkcolor=emerald, urlcolor=emerald, bookmarks=true]{hyperref}

\definecolor{emerald}{RGB}{0,102,51}
\definecolor{darkred}{RGB}{139,0,0}
\definecolor{lightgray}{RGB}{245,245,245}

\usepackage[nil]{babel}
\babelprovide[import, main]{bengali}
\babelprovide[import]{arabic}
\babelfont{rm}[Renderer=HarfBuzz,Scale=1.0]{Noto Serif Bengali}
\babelfont{sf}[Renderer=HarfBuzz]{Noto Sans Bengali}
\babelfont{tt}[Renderer=HarfBuzz]{Noto Sans Bengali}
\babelfont[arabic]{rm}[Renderer=HarfBuzz,Scale=1.1]{Noto Naskh Arabic}

% Section styling
\titleformat{\section}{\Large\bfseries\color{emerald}}{\thesection.}{0.5em}{}
\titleformat{\subsection}{\large\bfseries\color{emerald!85!black}}{\thesubsection.}{0.5em}{}
\titleformat{\subsubsection}{\normalsize\bfseries\color{emerald!70!black}}{\thesubsubsection.}{0.5em}{}
\titlespacing*{\section}{0pt}{18pt}{8pt}

% Fancy Header/Footer setup
\pagestyle{fancy}
\fancyhf{} % clear all header and footer fields
\fancyhead[L]{\color{emerald}\small\textbf{\nouppercase{\leftmark}}}
\fancyfoot[L]{\scriptsize\color{black!50}{\lat{AI}}-সংকলিত, আলেম-যাচাইকৃত নয়}
\fancyfoot[C]{\color{emerald!70!black}\thepage}
\renewcommand{\headrulewidth}{0.4pt}
\renewcommand{\headrule}{\hbox to\headwidth{\color{emerald}\leaders\hrule height \headrulewidth\hfill}}
\renewcommand{\footrulewidth}{0pt}

% Custom boxes
% 1. Ayat/Hadith Box
\newtcolorbox{ayatbox}[1][]{
    enhanced, breakable, 
    colback=emerald!5, 
    colframe=emerald, 
    boxrule=0.5pt, 
    arc=3pt, 
    left=10pt, right=10pt, top=10pt, bottom=10pt,
    #1
}

% 2. Quote/Translation Box
\newtcolorbox{quotebox}[1][]{
    enhanced, breakable,
    frame hidden,
    colback=lightgray,
    borderline west={3pt}{0pt}{emerald},
    left=15pt, right=10pt, top=10pt, bottom=10pt,
    sharp corners,
    #1
}

% 3. AI-generated notice box
\newtcolorbox{warnbox}[1][]{
    enhanced, breakable,
    colback=red!4,
    colframe=darkred,
    boxrule=0.8pt,
    arc=3pt,
    left=10pt, right=10pt, top=10pt, bottom=10pt,
    #1
}

% Commands
\newcommand{\arab}[1]{\foreignlanguage{arabic}{#1}}
\newcommand{\kib}[1]{\textcolor{emerald}{\textbf{#1}}}

% Latin (English) fallback: Noto Serif Bengali has NO Latin glyphs
\newfontfamily\latfont[Renderer=HarfBuzz]{Noto Serif}
\newcommand{\lat}[1]{{\latfont #1}}

\setlength{\parskip}{5pt}
\setlength{\parindent}{0pt}

% Fix for missing bullet in Noto Serif Bengali
\renewcommand{\labelitemi}{$\bullet$}



\begin{document}
\begin{center}{\Large\bfseries\color{emerald} হাদিস ও আসার — সালাত (নামাজ)}\end{center}
\vspace{1em}
\section*{হাদিস ও আসার — সালাত (নামাজ)}

\begin{warnbox}
 \textbf{গুরুত্বপূর্ণ নোটিশ (\lat{Important} \lat{Notice}):} এই কন্টেন্টটি \lat{AI} (কৃত্রিম বুদ্ধিমত্তা) দিয়ে প্রস্তুত ও সংকলিত। \lat{AI} কঠোর সূত্র-মান নীতি মেনে শুধু নির্ভরযোগ্য উৎস থেকে তথ্য সংগ্রহ করেছে — কেবল সহীহ/হাসান হাদিস ও সহীহ সনদের আসার রাখা হয়েছে, দুর্বল সনদ বাদ দেওয়া হয়েছে। তবে এটি কোনো আলেম বা মুহাদ্দিস কর্তৃক যাচাইকৃত নয়।
\end{warnbox}


\begin{quote}
\textbf{পড়ার নিয়ম:} প্রতিটি হাদিসের সাথে থাকবে — \textbf{সূত্র কার্ড} (সনদের মান) $\rightarrow$ \textbf{পটভূমি} $\rightarrow$ \textbf{পূর্ণ বর্ণনা (বাংলা)} $\rightarrow$ \textbf{বিশ্লেষণ} $\rightarrow$ \textbf{শিক্ষা}। \\
\textbf{সনদের মান:} বুখারি-মুসলিমের হাদিস \lat{“}সহীহ\lat{”}; অন্যান্য গ্রন্থের ক্ষেত্রে আল-আলবানী/তিরমিযীর শ্রেণি উল্লেখ করা হয়েছে। যে বর্ণনার সনদে মতভেদ আছে, তা স্পষ্টভাবে লেখা হয়েছে। ইবনে মাজাহ ১৪২৫-এর মতো দুর্বল সনদের সমান্তরাল বর্ণনা থাকলে তা লেবেলসহ দেখানো হয়েছে।
\end{quote}


\medskip\hrule\medskip

\subsection*{পার্ট ১ — নামাজের মর্যাদা ও গুরুত্ব}

\subsubsection*{হাদিস ১ — প্রথম হিসাব: নামাজ (সহীহ তিরমিযী ৪১৩)}

\#\#\#\# ১. সূত্র কার্ড

\begin{table}[h]\centering\small
\begin{tabular}{ll}
বিষয় & বিবরণ \\
\hline
\textbf{বর্ণনাকারী} & আবু হুরাইরা (রাঃ) \\
\textbf{গ্রন্থ} & সুনানে তিরমিযী ৪১৩; সুনানে আবু দাউদ ৮৬৪; সুনানে নাসাঈ ৪৬৫ \\
\textbf{অধ্যায়} & কিতাবুস-সালাত — বাবু মা জাআ ফি হিফজিস-সালাত \\
\textbf{সনদের মান} & \textbf{সহীহ} (আল-আলবানী, দারুসসালাম) \\
\textbf{দ্রষ্টব্য} & ইবনে মাজাহ ১৪২৫-এ একই হাদিস এসেছে, তবে সে সনদ দুর্বল (যয়ীফ) — এখানে সহীহ সনদগুলোই দলিল \\
\hline
\end{tabular}
\end{table}

\#\#\#\# ২. পটভূমি (\lat{Context})

নামাজ নিয়ে ইসলামের সবচেয়ে গুরুত্বপূর্ণ ঘোষণাগুলোর একটি — কিয়ামতের দিন \textbf{প্রথম হিসাব (\lat{first} \lat{reckoning})} নামাজের। এই হাদিসটি নামাজের মর্যাদা ও নামাজ ত্যাগ/ত্রুটির পরিণতি বোঝার মূল ভিত্তি।

\#\#\#\# ৩. পূর্ণ বর্ণনা (বাংলা, \lat{pure})

\begin{quote}
রাসূলুল্লাহ (সাল্লাল্লাহু আলাইহি ওয়া সাল্লাম) বলেছেন: \\
\textbf{\lat{“}কিয়ামতের দিন বান্দার আমলের মধ্যে সর্বপ্রথম নামাজের হিসাব নেওয়া হবে। যদি তা সঠিক হয়, তবে সে সফল ও সার্থক হবে; আর যদি তা নষ্ট হয়, তবে সে ব্যর্থ ও ক্ষতিগ্রস্ত হবে। আর ফরজ নামাজে ঘাটতি থাকলে আল্লাহ বলবেন — দেখো, আমার বান্দার কোনো নফল নামাজ আছে কি না, তা দিয়ে ফরজের ঘাটতি পূরণ করা হবে। তারপর বাকি আমলগুলোও এভাবেই হিসাব করা হবে।\lat{”}}
\end{quote}

\#\#\#\# ৪. বিশ্লেষণ (\lat{Analysis})

\textbf{১.} হিসাবের ক্রম — সর্বপ্রথম নামাজ: অর্থাৎ নামাজের \textbf{\lat{position} (অবস্থান)} ঈমানের পরই।
\textbf{২.} নফল দিয়ে ফরজের ঘাটতি পূরণ — নফল নামাজের গুরুত্ব এবং নামাজের ত্রুটির ক্ষতিপূরণের ব্যবস্থা।
\textbf{৩.} \lat{“}সঠিক/নষ্ট\lat{”} — এখানে সময়, পূর্ণতা ও মান সবই অন্তর্ভুক্ত।

\#\#\#\# ৫. শিক্ষা (\lat{Lessons})

\begin{enumerate}
\item নামাজের মান নিয়ে সচেতন থাকুন — এটি প্রথমে জিজ্ঞাসিত হবে।
\item ফরজের পাশাপাশি নফল (সুন্নত) নামাজ জমা করুন — ঘাটতি পূরণের \textbf{\lat{insurance} (নিরাপত্তা)}।
\end{enumerate}
\medskip\hrule\medskip

\subsubsection*{হাদিস ২ — নামাজ ত্যাগ: কুফরের সীমা (সহীহ মুসলিম ৮২)}

\#\#\#\# ১. সূত্র কার্ড

\begin{table}[h]\centering\small
\begin{tabular}{ll}
বিষয় & বিবরণ \\
\hline
\textbf{বর্ণনাকারী} & জাবির ইবনে আবদিল্লাহ (রাঃ) \\
\textbf{গ্রন্থ} & সহীহ মুসলিম ৮২; সুনানে তিরমিযী ২৬২০; সুনানে ইবনে মাজাহ ১০৭৮ \\
\textbf{অধ্যায়} & কিতাবুল ইমান — বাবু বায়ানিত তাসমিয়াহ \\
\textbf{সনদের মান} & \textbf{সহীহ} (বুখারি-মুসলিমের মানদণ্ডে, মুসলিম) \\
\hline
\end{tabular}
\end{table}

\#\#\#\# ২. পটভূমি (\lat{Context})

নামাজ ত্যাগের বিষয়টি ঈমানের সাথে যুক্ত — নবীজি (সা.) স্পষ্ট করে দিয়েছেন যে নামাজ ত্যাগ করা মানুষকে শিরক-কুফরের সীমায় পৌঁছে দেয়। এটি নামাজ ত্যাগের ভয়াবহতার সবচেয়ে শক্তিশালী দলিল।

\#\#\#\# ৩. পূর্ণ বর্ণনা (বাংলা, \lat{pure})

\begin{quote}
নবীজি (সাল্লাল্লাহু আলাইহি ওয়া সাল্লাম) বলেছেন: \\
\textbf{\lat{“}এক ব্যক্তি ও শিরক ও কুফরের মধ্যে পার্থক্য হলো নামাজ ত্যাগ করা।\lat{”}}
\end{quote}

\#\#\#\# ৪. বিশ্লেষণ (\lat{Analysis})

\textbf{১.} এই হাদিস নিয়ে আলেমদের দুটি ব্যাখ্যা: (ক) নামাজ ত্যাগ করা \textbf{কুফর} (সালাফের অনেকের মত — ইবনে হাযম, ইবনে তাইমিয়্যার প্রসিদ্ধ মত; কারণ হাদিসের বাহ্যিক অর্থ), (খ) ত্যাগকারীর আমল কুফরের \textbf{আমল} হয়, কিন্তু সে মুসলিম থাকে — যতক্ষণ না নামাজ ত্যাগকে বৈধ মনে করে (জমহুর ফকীহগণ)। ব্যাখ্যার পার্থক্য \textbf{নামাজ ত্যাগের গুরুত্ব} সম্পর্কে নয় — সবাই একমত এটি সর্বোচ্চ পর্যায়ের অন্যায়।
\textbf{২.} হাদিসটির সমর্থনে তিরমিযী ২৬২১-এর হাদিস — \lat{“}আমাদের ও তাদের মধ্যে অঙ্গীকার হলো নামাজ; যে তা ত্যাগ করল, সে কুফরি করল\lat{”} (সহীহ)।

\#\#\#\# ৫. শিক্ষা (\lat{Lessons})

\begin{enumerate}
\item নামাজ ত্যাগকে হালকা ভাবার কোনো সুযোগ নেই — এটি ঈমানের সীমানার বিষয়।
\item মতভেদ থাকলেও সবার ঐকমত্য — নামাজ ত্যাগ \textbf{জঘন্য পাপ}; দ্রুত তাওবায় ফিরে আসা ফরজ।
\end{enumerate}
\medskip\hrule\medskip

\subsubsection*{হাদিস ৩ — অঙ্গীকার: নামাজ ত্যাগকারীর চুক্তিভঙ্গ (সহীহ তিরমিযী ২৬২১)}

\#\#\#\# ১. সূত্র কার্ড

\begin{table}[h]\centering\small
\begin{tabular}{ll}
বিষয় & বিবরণ \\
\hline
\textbf{বর্ণনাকারী} & বুরাইদা ইবনুল হুসাইব (রাঃ) \\
\textbf{গ্রন্থ} & সুনানে তিরমিযী ২৬২১; সুনানে ইবনে মাজাহ ১০৭৯; সুনানে নাসাঈ ৪৬৩ \\
\textbf{অধ্যায়} & কিতাবুল ইমান — বাবু মা জাআ ফি তারকিস-সালাত \\
\textbf{সনদের মান} & \textbf{সহীহ} (আল-আলবানী) \\
\hline
\end{tabular}
\end{table}

\#\#\#\# ২. পটভূমি (\lat{Context})

নবীজি (সা.) মুসলিম ও কাফিরের মধ্যকার চুক্তির কথা বলেছেন — সেই চুক্তির কেন্দ্রবিন্দু নামাজ।

\#\#\#\# ৩. পূর্ণ বর্ণনা (বাংলা, \lat{pure})

\begin{quote}
\textbf{\lat{“}আমাদের এবং তাদের মধ্যে (প্রতিষ্ঠিত) অঙ্গীকার হলো নামাজ। যে ব্যক্তি তা ত্যাগ করল, সে (আসলে) কুফরি করল।\lat{”}}
\end{quote}

\#\#\#\# ৪. বিশ্লেষণ (\lat{Analysis})

\textbf{১.} \lat{“}অঙ্গীকার\lat{”} — ইসলাম গ্রহণের পর নামাজই ঈমানের বাস্তব \textbf{প্রকাশ (\lat{manifestation})}।
\textbf{২.} এই হাদিস হাদিস ২-এর সমর্থক — একত্রে নামাজ ত্যাগের সতর্কবার্তার ভিত্তি।

\#\#\#\# ৫. শিক্ষা (\lat{Lessons})

\begin{enumerate}
\item নামাজ = চুক্তির প্রতীক — প্রতিদিনের নামাজ এই চুক্তিকে জীবন্ত রাখে।
\item ত্যাগকারীকে ফিরিয়ে আনার দায়িত্ব পরিবারের \textbf{\lat{guardian} (অভিভাবক)}-এর।
\end{enumerate}
\medskip\hrule\medskip

\subsubsection*{হাদিস ৪ — দ্বীনের স্তম্ভ: নামাজ (সহীহ তিরমিযী ২৬১৬)}

\#\#\#\# ১. সূত্র কার্ড

\begin{table}[h]\centering\small
\begin{tabular}{ll}
বিষয় & বিবরণ \\
\hline
\textbf{বর্ণনাকারী} & মুআয ইবনে জাবাল (রাঃ) \\
\textbf{গ্রন্থ} & সুনানে তিরমিযী ২৬১৬; সুনানে ইবনে মাজাহ ২৮০৬ \\
\textbf{অধ্যায়} & কিতাবুল ইমান \\
\textbf{সনদের মান} & \textbf{হাসান সহীহ} (তিরমিযী); হাসান (আল-আলবানী) \\
\hline
\end{tabular}
\end{table}

\#\#\#\# ২. পটভূমি (\lat{Context})

মুআয (রাঃ) উটের পিঠে নবীজির সাথে ছিলেন — নবীজি তাকে দ্বীনের সারকথা শিখিয়েছেন।

\#\#\#\# ৩. পূর্ণ বর্ণনা (বাংলা, \lat{pure})

\begin{quote}
\textbf{\lat{“}আমি কি তোমাকে দ্বীনের মূল কথা বলব না? দ্বীনের মাথা হলো ইসলাম, তার স্তম্ভ হলো নামাজ, আর তার সর্বোচ্চ শিখর হলো জ\lat{ihad} (জিহাদ)।\lat{”}}
\end{quote}

\#\#\#\# ৪. বিশ্লেষণ (\lat{Analysis})

\textbf{১.} \lat{“}স্তম্ভ\lat{”} — ঘর যেমন স্তম্ভ ছাড়া টেকে না, দ্বীনও নামাজ ছাড়া \textbf{\lat{standing} (টিকে থাকা)} সম্ভব নয় — এই রূপকই নামাজের মর্যাদা বুঝায়।
\textbf{২.} জনপ্রিয় বাক্য \lat{“}নামাজ দ্বীনের স্তম্ভ\lat{”} — এর নির্ভরযোগ্য উৎস এই হাদিস; আলাদা কোনো সহীহ হাদিসে \lat{“}নামাজ দ্বীনের স্তম্ভ\lat{”} শব্দে আসেনি।

\#\#\#\# ৫. শিক্ষা (\lat{Lessons})

\begin{enumerate}
\item নামাজ দ্বীনের কাঠামোর \textbf{\lat{core} (মূল)} — একে দুর্বল করলে পুরো কাঠামো ঝুঁকিতে।
\item এই হাদিসটি মনে রাখুন — তর্ক-বিতর্কে দলিল হিসেবে ব্যবহারযোগ্য।
\end{enumerate}
\medskip\hrule\medskip

\subsubsection*{হাদিস ৫ — আসর ছুটলে: পরিবার-ধন হারানো (সহীহ বুখারি ৫৫২)}

\#\#\#\# ১. সূত্র কার্ড

\begin{table}[h]\centering\small
\begin{tabular}{ll}
বিষয় & বিবরণ \\
\hline
\textbf{বর্ণনাকারী} & আবদুল্লাহ ইবনে উমার (রাঃ) \\
\textbf{গ্রন্থ} & সহীহ বুখারি ৫৫২; সহীহ মুসলিম ৬২৬ \\
\textbf{অধ্যায়} & কিতাবুল মাওয়াকীত — বাবু ইসমিত তাসমিয়াহ \\
\textbf{সনদের মান} & \textbf{সহীহ} (বুখারি-মুসলিম) \\
\hline
\end{tabular}
\end{table}

\#\#\#\# ২. পটভূমি (\lat{Context})

আসরের নামাজ — দিনের ব্যস্ততার শেষে; নবীজি (সা.) আসর ছুটে যাওয়াকে সবচেয়ে বড় ক্ষতির সাথে তুলনা করেছেন।

\#\#\#\# ৩. পূর্ণ বর্ণনা (বাংলা, \lat{pure})

\begin{quote}
\textbf{\lat{“}যে ব্যক্তি আসরের নামাজ ছেড়ে দিল, সে যেন তার পরিবার-পরিজন ও ধন-সম্পদ হারিয়ে ফেলল।\lat{”}}
\end{quote}

\#\#\#\# ৪. বিশ্লেষণ (\lat{Analysis})

\textbf{১.} রূপকটি অত্যন্ত শক্তিশালী — পরিবার ও সম্পদ মানুষের কাছে সবচেয়ে প্রিয়; আসর ছুটে যাওয়া তার চেয়েও বড় \textbf{\lat{loss} (ক্ষতি)}।
\textbf{২.} এই হাদিস সূরা বাকারাহ ২:২৩৮-এর \lat{“}সালাতুল উস্তা\lat{”} (আসর) রক্ষার নির্দেশের সাথে মিলে যায়।

\#\#\#\# ৫. শিক্ষা (\lat{Lessons})

\begin{enumerate}
\item আসরের নামাজের জন্য বিশেষ \textbf{\lat{alertness} (সতর্কতা)} — কাজের চাপে যেন না ছুটে।
\item প্রতিটি নামাজকে পরিবার-ধনের মতোই মূল্যবান মনে করুন।
\end{enumerate}
\medskip\hrule\medskip

\subsubsection*{হাদিস ৬ — মুনাফিকদের উপর ভারী নামাজ: ইশা ও ফজর (সহীহ বুখারি ৬৫৭)}

\#\#\#\# ১. সূত্র কার্ড

\begin{table}[h]\centering\small
\begin{tabular}{ll}
বিষয় & বিবরণ \\
\hline
\textbf{বর্ণনাকারী} & আবু হুরাইরা (রাঃ) \\
\textbf{গ্রন্থ} & সহীহ বুখারি ৬৫৭; সহীহ মুসলিম ৬৫১ \\
\textbf{অধ্যায়} & কিতাবুল আজান — বাবু ফজলিল ইশা \\
\textbf{সনদের মান} & \textbf{সহীহ} (বুখারি-মুসলিম) \\
\hline
\end{tabular}
\end{table}

\#\#\#\# ২. পটভূমি (\lat{Context})

মুনাফিকদের নামাজে অনীহার কারণ বিশ্লেষণ করতে গিয়ে নবীজি (সা.) ইশা ও ফজরের জামাতের কথা বলেছেন — অন্ধকার ও ঘুমের কারণে এই দুই নামাজে জামাত মুনাফিকদের জন্য ভারী।

\#\#\#\# ৩. পূর্ণ বর্ণনা (বাংলা, \lat{pure})

\begin{quote}
\textbf{\lat{“}মুনাফিকদের উপর সবচেয়ে ভারী নামাজ হলো ইশা ও ফজরের নামাজ। তারা যদি জানত, এই দুই নামাজে (জামাতে) কী সওয়াব আছে, তবে হামাগুড়ি দিয়ে হলেও (মসজিদে) আসত।\lat{”}}
\end{quote}

\#\#\#\# ৪. বিশ্লেষণ (\lat{Analysis})

\textbf{১.} জামাতের মূল্য — এই দুই নামাজের জামাতই মুনাফিকি ও ঈমানের \textbf{পরীক্ষা (\lat{test})}।
\textbf{২.} ফজর-ইশার জামাতে উপস্থিতি ঈমানের শক্তি ও নামাজের ভালোবাসার পরিচয়।

\#\#\#\# ৫. শিক্ষা (\lat{Lessons})

\begin{enumerate}
\item ফজর ও ইশা — জামাতের জন্য সর্বোচ্চ চেষ্টা করুন।
\item অলসতা কাটিয়ে ওঠার জন্য রাতে তাড়াতাড়ি ঘুমানোর \textbf{\lat{routine} (অভ্যাস)}।
\end{enumerate}
\medskip\hrule\medskip

\subsubsection*{হাদিস ৭ — পাঁচ নামাজ: গুনাহের কাফফারা (সহীহ মুসলিম ২৩৩)}

\#\#\#\# ১. সূত্র কার্ড

\begin{table}[h]\centering\small
\begin{tabular}{ll}
বিষয় & বিবরণ \\
\hline
\textbf{বর্ণনাকারী} & আবু হুরাইরা (রাঃ) \\
\textbf{গ্রন্থ} & সহীহ মুসলিম ২৩৩ \\
\textbf{অধ্যায়} & কিতাবুত তাহারাহ — বাবু মাশরুইয়াতি মুলাযামাতিল কুলূব \\
\textbf{সনদের মান} & \textbf{সহীহ} \\
\hline
\end{tabular}
\end{table}

\#\#\#\# ২. পটভূমি (\lat{Context})

পাঁচ ওয়াক্ত নামাজ ও জুমা — গুনাহ মোচনের আল্লাহর দেওয়া নিয়মিত সুযোগ।

\#\#\#\# ৩. পূর্ণ বর্ণনা (বাংলা, \lat{pure})

\begin{quote}
\textbf{\lat{“}পাঁচ ওয়াক্ত নামাজ এবং জুমা থেকে পরবর্তী জুমা — এগুলোর মধ্যবর্তী গুনাহের কাফফারা, যদি বড় গুনাহ (কবিরাহ) না করা হয়।\lat{”}}
\end{quote}

\#\#\#\# ৪. বিশ্লেষণ (\lat{Analysis})

\textbf{১.} শর্ত — বড় গুনাহ থেকে বিরত থাকা; নামাজ ছোট গুনাহ মুছে দেয়, বড় গুনাহের জন্য তাওবা প্রয়োজন।
\textbf{২.} নামাজের \textbf{\lat{purifying} (পবিত্রকরণ)} ভূমিকা — দৈনন্দিন পাপ ধুয়ে যায় পাঁচবারের নামাজে।

\#\#\#\# ৫. শিক্ষা (\lat{Lessons})

\begin{enumerate}
\item পাপের বোঝা নিয়ে হতাশ হবেন না — নামাজই আল্লাহর ক্ষমার দরজা।
\item তবে বড় গুনাহে লিপ্ত হলে নামাজের পাশাপাশি খাঁটি তাওবা করুন।
\end{enumerate}
\medskip\hrule\medskip

\subsubsection*{হাদিস ৮ — চোখের শীতলতা: নামাজ (সহীহ নাসাঈ ৩৯৩৯)}

\#\#\#\# ১. সূত্র কার্ড

\begin{table}[h]\centering\small
\begin{tabular}{ll}
বিষয় & বিবরণ \\
\hline
\textbf{বর্ণনাকারী} & আনাস ইবনে মালিক (রাঃ) \\
\textbf{গ্রন্থ} & সুনানে নাসাঈ ৩৯৩৯; মুসনাদ আহমাদ \\
\textbf{অধ্যায়} & কিতাবু ইশরাতিন নিসা — বাবু হুব্বিন নিসা \\
\textbf{সনদের মান} & \textbf{হাসান} (দারুসসালাম; আল-আলবানী) \\
\hline
\end{tabular}
\end{table}

\#\#\#\# ২. পটভূমি (\lat{Context})

নবীজি (সা.) দুনিয়ার তিনটি জিনিস প্রিয় বলেছেন — নারী, সুগন্ধি; আর চোখের শীতলতা নামাজে।

\#\#\#\# ৩. পূর্ণ বর্ণনা (বাংলা, \lat{pure})

\begin{quote}
\textbf{\lat{“}আমার কাছে দুনিয়ার তিনটি জিনিস প্রিয় করা হয়েছে — নারী ও সুগন্ধি; আর চোখের শীতলতা রাখা হয়েছে নামাজে।\lat{”}}
\end{quote}

\#\#\#\# ৪. বিশ্লেষণ (\lat{Analysis})

\textbf{১.} \lat{“}চোখের শীতলতা\lat{”} — আরবি অভিব্যক্তি: সবচেয়ে বড় আনন্দ ও প্রশান্তি; নবীজির কাছে নামাজই সেই \textbf{\lat{tranquility} (প্রশান্তি)}।
\textbf{২.} নামাজ নবীজির জন্য বোঝা নয়, বিশ্রাম — এই দৃষ্টিভঙ্গিই নামাজের প্রকৃত ভালোবাসা।

\#\#\#\# ৫. শিক্ষা (\lat{Lessons})

\begin{enumerate}
\item নামাজকে ভালোবাসা শিখুন — বোঝা ভাবলে খুশু আসে না।
\item নামাজে প্রশান্তি খুঁজুন — দুশ্চিন্তার ওষুধ নামাজ।
\end{enumerate}
\medskip\hrule\medskip

\subsubsection*{হাদিস ৯ — শেষ অসিয়ত: নামাজ, নামাজ (সহীহ আবু দাউদ ৫১৫৬)}

\#\#\#\# ১. সূত্র কার্ড

\begin{table}[h]\centering\small
\begin{tabular}{ll}
বিষয় & বিবরণ \\
\hline
\textbf{বর্ণনাকারী} & উম্মে সালামাহ (রাঃ) \\
\textbf{গ্রন্থ} & সুনানে আবু দাউদ ৫১৫৬; সুনানে ইবনে মাজাহ ২৬৯৮ \\
\textbf{অধ্যায়} & কিতাবুল আদাব — বাবু ফি ইসতিসওয়ারিল আবদ \\
\textbf{সনদের মান} & \textbf{সহীহ} (আল-আলবানী) \\
\hline
\end{tabular}
\end{table}

\#\#\#\# ২. পটভূমি (\lat{Context})

মৃত্যুশয্যায় নবীজি (সা.)-এর শেষ কথাগুলোর মধ্যে নামাজ ছিল — তার মুখে বারবার ফিরে আসছিল: \lat{“}নামাজ, নামাজ...\lat{”}

\#\#\#\# ৩. পূর্ণ বর্ণনা (বাংলা, \lat{pure})

\begin{quote}
উম্মে সালামাহ (রাঃ) বলেন — নবীজি (সাল্লাল্লাহু আলাইহি ওয়া সাল্লাম) মৃত্যুর আগে শেষ যে কথাগুলো বলেছিলেন: \\
\textbf{\lat{“}নামাজ, নামাজ! আর তোমাদের অধীনস্থদের (দাস-দাসীদের) ব্যাপারে আল্লাহকে ভয় করো।\lat{”}} — তিনি এ কথা বারবার বলতে থাকলেন, যতক্ষণ না তাঁর জিহবা স্তব্ধ হয়ে যায়।
\end{quote}

\#\#\#\# ৪. বিশ্লেষণ (\lat{Analysis})

\textbf{১.} মৃত্যুশয্যার অসিয়ত — জীবনের সবচেয়ে গুরুত্বপূর্ণ বার্তা; নবীজির কাছে সেটিই নামাজ।
\textbf{২.} অসিয়তের \textbf{\lat{sequence} (ক্রম)} — নামাজ আগে, তারপর মানুষের অধিকার।

\#\#\#\# ৫. শিক্ষা (\lat{Lessons})

\begin{enumerate}
\item নিজের \lat{“}শেষ কথা\lat{”} কল্পনা করুন — নামাজ যেন সেই তালিকায় প্রথম থাকে।
\item অধীনস্থ (কর্মচারী, সন্তান, গৃহকর্মী) — তাদের অধিকার রক্ষা নামাজের পাশাপাশি।
\end{enumerate}
\medskip\hrule\medskip

\subsubsection*{হাদিস ১০ — ফজরের দুই রাকাত: দুনিয়ার চেয়ে উত্তম (সহীহ মুসলিম ৭২৫)}

\#\#\#\# ১. সূত্র কার্ড

\begin{table}[h]\centering\small
\begin{tabular}{ll}
বিষয় & বিবরণ \\
\hline
\textbf{বর্ণনাকারী} & আয়িশা (রাঃ) \\
\textbf{গ্রন্থ} & সহীহ মুসলিম ৭২৫ \\
\textbf{অধ্যায়} & কিতাবুস-সালাত — বাবু রাফিল ইয়াদাইন \\
\textbf{সনদের মান} & \textbf{সহীহ} \\
\hline
\end{tabular}
\end{table}

\#\#\#\# ২. পটভূমি (\lat{Context})

ফজরের দুই রাকাত সুন্নত — নবীজি (সা.) কখনো তা ছাড়তেন না; এখানে তার মর্যাদা ঘোষিত।

\#\#\#\# ৩. পূর্ণ বর্ণনা (বাংলা, \lat{pure})

\begin{quote}
\textbf{\lat{“}ফজরের (ফরজের আগের) দুই রাকাত (সুন্নত) দুনিয়া ও দুনিয়ার যা কিছু আছে — সবকিছু থেকে উত্তম।\lat{”}}
\end{quote}

\#\#\#\# ৪. বিশ্লেষণ (\lat{Analysis})

\textbf{১.} দুই রাকাতের এত মর্যাদা — কারণ ভোরের ইবাদতে অন্তর সবচেয়ে \textbf{সজাগ (\lat{awake})} থাকে।
\textbf{২.} নবীজি ফজরের সুন্নত কখনো ছাড়েননি — এমনকি সফরেও (মুসলিম ৭২৫)।

\#\#\#\# ৫. শিক্ষা (\lat{Lessons})

\begin{enumerate}
\item ফজরের দুই রাকাত সুন্নতকে অপরিহার্য ভাবুন — ছুটলে সূর্যোদয়ের পর কাযা করা যায় (তিরমিযী ৪২৩)।
\item ভোরে ওঠার \textbf{\lat{habit} (অভ্যাস)} গড়ুন — ফজরের সুন্নতই দিনের প্রথম বিজয়।
\end{enumerate}
\medskip\hrule\medskip

\subsubsection*{হাদিস ১১ — জামাতে ফজর-ইশা: সারা রাতের ইবাদত (সহীহ মুসলিম ৬৫৬)}

\#\#\#\# ১. সূত্র কার্ড

\begin{table}[h]\centering\small
\begin{tabular}{ll}
বিষয় & বিবরণ \\
\hline
\textbf{বর্ণনাকারী} & উসমান ইবনে আফফান (রাঃ) \\
\textbf{গ্রন্থ} & সহীহ মুসলিম ৬৫৬ \\
\textbf{অধ্যায়} & কিতাবুস-সালাত — বাবু ফজলি সালাতিল জামাআহ \\
\textbf{সনদের মান} & \textbf{সহীহ} \\
\hline
\end{tabular}
\end{table}

\#\#\#\# ২. পটভূমি (\lat{Context})

জামাতের পুরস্কারের তুলনা — নবীজি (সা.) ইশা ও ফজরের জামাতকে রাত জাগরণের সমান বলেছেন।

\#\#\#\# ৩. পূর্ণ বর্ণনা (বাংলা, \lat{pure})

\begin{quote}
\textbf{\lat{“}যে ব্যক্তি জামাতে ইশার নামাজ পড়ল, সে যেন অর্ধরাত্রি ইবাদতে কাটাল; আর যে ব্যক্তি জামাতে ফজরের নামাজ পড়ল, সে যেন সমস্ত রাত ইবাদতে কাটাল।\lat{”}}
\end{quote}

\#\#\#\# ৪. বিশ্লেষণ (\lat{Analysis})

\textbf{১.} দুই নামাজ মিলে সম্পূর্ণ রাতের ইবাদতের সওয়াব — অল্প সময়ে বিশাল \textbf{\lat{reward} (পুরস্কার)}।
\textbf{২.} ফজর-ইশার জামাত নিয়মিত করলে রাতের তাহাজ্জুদের সওয়াবের দরজা খুলে যায়।

\#\#\#\# ৫. শিক্ষা (\lat{Lessons})

\begin{enumerate}
\item ইশা ও ফজর — জামাতের দুটি সোনার সুযোগ।
\item যারা তাহাজ্জুদ পারে না, তারা এই দুই নামাজের জামাতে \textbf{\lat{compensate} (ক্ষতিপূরণ)} করুন।
\end{enumerate}
\medskip\hrule\medskip

\subsection*{পার্ট ২ — নামাজের ভিত্তি, শর্ত ও জামাত}

\subsubsection*{হাদিস ১২ — নামাজের চাবি: পবিত্রতা (সহীহ আবু দাউদ ৬১)}

\#\#\#\# ১. সূত্র কার্ড

\begin{table}[h]\centering\small
\begin{tabular}{ll}
বিষয় & বিবরণ \\
\hline
\textbf{বর্ণনাকারী} & আলী ইবনে আবি তালিব (রাঃ) \\
\textbf{গ্রন্থ} & সুনানে আবু দাউদ ৬১; সুনানে তিরমিযী ৪; সুনানে ইবনে মাজাহ ২৭৫ \\
\textbf{অধ্যায়} & কিতাবুত তাহারাহ — বাবু মাফতাহিস-সালাত \\
\textbf{সনদের মান} & \textbf{সহীহ} (আল-আলবানী) \\
\hline
\end{tabular}
\end{table}

\#\#\#\# ২. পটভূমি (\lat{Context})

নামাজের প্রবেশদ্বার ও প্রস্থানদ্বার — নবীজি (সা.) এক বাক্যে নামাজের সীমানা ঘোষণা করেছেন।

\#\#\#\# ৩. পূর্ণ বর্ণনা (বাংলা, \lat{pure})

\begin{quote}
\textbf{\lat{“}নামাজের চাবি হলো পবিত্রতা (তাহারাত); নামাজে প্রবেশকারী নিষিদ্ধকারী হলো তাকবীর (তাকবীরে তাহরিমা); আর নামাজ থেকে বেরকারী হলো সালাম।\lat{”}}
\end{quote}

\#\#\#\# ৪. বিশ্লেষণ (\lat{Analysis})

\textbf{১.} তিনটি সীমানা — (ক) পবিত্রতা ছাড়া নামাজ শুরুই হয় না, (খ) তাকবীরের পর দুনিয়ার কাজ হারাম, (গ) সালামের মাধ্যমে নামাজ শেষ।
\textbf{২.} তাকবীরে তাহরিমার গুরুত্ব — নামাজের \textbf{\lat{entrance} (প্রবেশ)} এখানেই; কোনো কাজ যদি তাকবীরের আগে হয়, তা নামাজের অংশ নয়।

\#\#\#\# ৫. শিক্ষা (\lat{Lessons})

\begin{enumerate}
\item ওযুর মান যাচাই করুন — নামাজের চাবিই পবিত্রতা।
\item তাকবীর থেকে সালাম পর্যন্ত — নামাজের \textbf{\lat{boundary} (সীমানা)} মেনে চলুন।
\end{enumerate}
\medskip\hrule\medskip

\subsubsection*{হাদিস ১৩ — পবিত্রতা ছাড়া নামাজ কবুল নয় (সহীহ মুসলিম ২২৪)}

\#\#\#\# ১. সূত্র কার্ড

\begin{table}[h]\centering\small
\begin{tabular}{ll}
বিষয় & বিবরণ \\
\hline
\textbf{বর্ণনাকারী} & আবদুল্লাহ ইবনে উমার (রাঃ) \\
\textbf{গ্রন্থ} & সহীহ মুসলিম ২২৪; সুনানে তিরমিযী ১; সুনানে ইবনে মাজাহ ২৭১ \\
\textbf{অধ্যায়} & কিতাবুত তাহারাহ — বাবু উজুবিত তাহারাহ \\
\textbf{সনদের মান} & \textbf{সহীহ} \\
\hline
\end{tabular}
\end{table}

\#\#\#\# ২. পটভূমি (\lat{Context})

ওযু ছাড়া নামাজের \textbf{\lat{invalidity} (বাতিলতা)} — ইসলামের সবচেয়ে পরিচিত বিধানগুলোর ভিত্তি।

\#\#\#\# ৩. পূর্ণ বর্ণনা (বাংলা, \lat{pure})

\begin{quote}
\textbf{\lat{“}পবিত্রতা (ওযু) ছাড়া আল্লাহ নামাজ কবুল করেন না; আর (জায়েজ সম্পদ থেকে) খিয়ানতের মাল দিয়ে দান-সদকাও কবুল হয় না।\lat{”}}
\end{quote}

\#\#\#\# ৪. বিশ্লেষণ (\lat{Analysis})

\textbf{১.} নামাজের শর্ত — ওযু/পবিত্রতা; শর্তহীন নামাজ আদায়ই হয়নি (বাতিল)।
\textbf{২.} হাদিসের দ্বিতীয় অংশ — উপার্জনের পবিত্রতাও ইবাদত কবুলের শর্ত।

\#\#\#\# ৫. শিক্ষা (\lat{Lessons})

\begin{enumerate}
\item নামাজের আগে ওযু \textbf{নিয়মমতো (\lat{properly})} করুন — তাড়াহুড়া নয়।
\item হালাল উপার্জনের চেষ্টা করুন — কারণ এর সাথে ইবাদত কবুলের সম্পর্ক।
\end{enumerate}
\medskip\hrule\medskip

\subsubsection*{হাদিস ১৪ — যেমন দেখেছ আমাকে নামাজ পড়তে (সহীহ বুখারি ৬৩১)}

\#\#\#\# ১. সূত্র কার্ড

\begin{table}[h]\centering\small
\begin{tabular}{ll}
বিষয় & বিবরণ \\
\hline
\textbf{বর্ণনাকারী} & মালিক ইবনে হুওয়াইরিস (রাঃ) \\
\textbf{গ্রন্থ} & সহীহ বুখারি ৬৩১ \\
\textbf{অধ্যায়} & কিতাবুল আজান — বাবুল আজান লিমান ইয়ারিদু জামাআহ \\
\textbf{সনদের মান} & \textbf{সহীহ} \\
\hline
\end{tabular}
\end{table}

\#\#\#\# ২. পটভূমি (\lat{Context})

দুই যুবক সাহাবি নবীজির কাছে এসেছিলেন — তাদের বিদায়ের সময় নবীজি তাদের নামাজ শিখিয়ে বললেন — এটি নামাজের পদ্ধতির সবচেয়ে বড় ভিত্তি।

\#\#\#\# ৩. পূর্ণ বর্ণনা (বাংলা, \lat{pure})

\begin{quote}
মালিক ইবনে হুওয়াইরিস (রাঃ) বলেন — নবীজি (সাল্লাল্লাহু আলাইহি ওয়া সাল্লাম) আমাদের বলেছেন: \\
\textbf{\lat{“}তোমরা নামাজ পড়ো, যেমন দেখেছ আমাকে নামাজ পড়তে।\lat{”}}
\end{quote}

\#\#\#\# ৪. বিশ্লেষণ (\lat{Analysis})

\textbf{১.} নামাজের পদ্ধতির \textbf{\lat{authority} (উৎস)} — নবীজির অনুসরণ; এখানেই ফিকহের সব মাযহাবের মিল।
\textbf{২.} এই হাদিস নামাজের পদ্ধতিগত মতভেদ (রফউল ইয়াদাইন, হাত বাঁধা ইত্যাদি) মীমাংসার মানদণ্ড — প্রতিটি মাযহাবই দাবি করে তার পদ্ধতি নবীজির আমলের নিকটতম।

\#\#\#\# ৫. শিক্ষা (\lat{Lessons})

\begin{enumerate}
\item নামাজের পদ্ধতি শিখুন নবীজির সুন্নত থেকে — অভ্যাস থেকে নয়।
\item মতভেদের ক্ষেত্রে প্রশ্ন করুন — \lat{“}নবীজি কীভাবে পড়েছেন?\lat{”}
\end{enumerate}
\medskip\hrule\medskip

\subsubsection*{হাদিস ১৫ — ফাতিহা ছাড়া নামাজ নেই (সহীহ বুখারি ৭৫৬)}

\#\#\#\# ১. সূত্র কার্ড

\begin{table}[h]\centering\small
\begin{tabular}{ll}
বিষয় & বিবরণ \\
\hline
\textbf{বর্ণনাকারী} & উবাদা ইবনুস সামিত (রাঃ) \\
\textbf{গ্রন্থ} & সহীহ বুখারি ৭৫৬; সহীহ মুসলিম ৩৯৪ \\
\textbf{অধ্যায়} & কিতাবুল আজান — বাবু উজুবিল কিরাআতিল ইমাম \\
\textbf{সনদের মান} & \textbf{সহীহ} (বুখারি-মুসলিম) \\
\hline
\end{tabular}
\end{table}

\#\#\#\# ২. পটভূমি (\lat{Context})

সূরা ফাতিহা নামাজের অপরিহার্য অংশ — নবীজি (সা.) একে \lat{“}নামাজের ভিত্তি\lat{”} বলেছেন; ইমামের পেছনে ফাতিহা পড়ার মতভেদ এই হাদিসের ব্যাখ্যাকে ঘিরে (বিস্তারিত — পার্ট ৬-এর ফাইল)।

\#\#\#\# ৩. পূর্ণ বর্ণনা (বাংলা, \lat{pure})

\begin{quote}
\textbf{\lat{“}যে ব্যক্তি সূরা ফাতিহা পড়েনি, তার নামাজ নেই (পূর্ণ হয় না)।\lat{”}}
\end{quote}

\#\#\#\# ৪. বিশ্লেষণ (\lat{Analysis})

\textbf{১.} জমহুরের মতে — ফাতিহা প্রতিটি রাকাতের রুকন; ছাড়লে নামাজ বাতিল।
\textbf{২.} ইমামের পেছনে মুক্তাদির ফাতিহা — এ নিয়ে চার মাযহাবে মতভেদ; কিন্তু একাকী নামাজে ফাতিহা নিয়ে কোনো মতভেদ নেই — \textbf{সর্বসম্মত (\lat{consensus})}।

\#\#\#\# ৫. শিক্ষা (\lat{Lessons})

\begin{enumerate}
\item ফাতিহা শুদ্ধভাবে পড়া শিখুন — উচ্চারণ ঠিক রাখুন।
\item মুক্তাদি হলে মাযহাবের নিয়ম মেনে চলুন — তবে ইমামের নীরব কিরাআতে ফাতিহা পড়া অধিকাংশ আলেমের কাছে নিরাপদ (বিস্তারিত — পার্ট ৬)।
\end{enumerate}
\medskip\hrule\medskip

\subsubsection*{হাদিস ১৬ — জমিনকে মসজিদ করা হয়েছে (সহীহ বুখারি ৩৩৫)}

\#\#\#\# ১. সূত্র কার্ড

\begin{table}[h]\centering\small
\begin{tabular}{ll}
বিষয় & বিবরণ \\
\hline
\textbf{বর্ণনাকারী} & জাবির ইবনে আবদিল্লাহ (রাঃ) \\
\textbf{গ্রন্থ} & সহীহ বুখারি ৩৩৫; সহীহ মুসলিম ৫২১ \\
\textbf{অধ্যায়} & কিতাবুত তায়াম্মুম — বাবুল মসজিদ \\
\textbf{সনদের মান} & \textbf{সহীহ} \\
\hline
\end{tabular}
\end{table}

\#\#\#\# ২. পটভূমি (\lat{Context})

নবীজিকে দেওয়া পাঁচটি বিশেষত্বের মধ্যে একটি — গোটা জমিন নামাজের জায়গা ও তায়াম্মুমের উপকরণ।

\#\#\#\# ৩. পূর্ণ বর্ণনা (বাংলা, \lat{pure})

\begin{quote}
\textbf{\lat{“}আমাকে পাঁচটি জিনিস দেওয়া হয়েছে, যা আমার আগে কাউকে দেওয়া হয়নি: ... আমার জন্য জমিনকে মসজিদ ও পবিত্রতা (তায়াম্মুমের মাধ্যম) বানানো হয়েছে। তাই আমার উম্মতের কোনো ব্যক্তি নামাজের সময় পেলে (যেখানেই থাকুক) নামাজ পড়ে নেবে।\lat{”}}
\end{quote}

\#\#\#\# ৪. বিশ্লেষণ (\lat{Analysis})

\textbf{১.} মসজিদের প্রয়োজনীয়তা নয় — নামাজের \textbf{\lat{validity} (শুদ্ধতা)} জায়গার উপর নির্ভরশীল নয়; পরিচ্ছন্ন জায়গাই যথেষ্ট।
\textbf{২.} তায়াম্মুমের ভিত্তি — পানি না পেলে মাটি দিয়ে পবিত্রতার বিকল্প।

\#\#\#\# ৫. শিক্ষা (\lat{Lessons})

\begin{enumerate}
\item মসজিদ থেকে দূরে থাকলে নামাজ ছাড়ার অজুহাত নেই — যে পরিচ্ছন্ন জায়গাতেই পড়ুন।
\item সফরে নামাজের \textbf{\lat{facility} (সুবিধা)} — ইসলাম কঠোরতা নয়।
\end{enumerate}
\medskip\hrule\medskip

\subsubsection*{হাদিস ১৭ — জামাত: ২৫/২৭ গুণ মর্যাদা (সহীহ বুখারি ৬৪৫)}

\#\#\#\# ১. সূত্র কার্ড

\begin{table}[h]\centering\small
\begin{tabular}{ll}
বিষয় & বিবরণ \\
\hline
\textbf{বর্ণনাকারী} & আবদুল্লাহ ইবনে উমার (রাঃ) \\
\textbf{গ্রন্থ} & সহীহ বুখারি ৬৪৫; সহীহ মুসলিম ৬৫০ \\
\textbf{অধ্যায়} & কিতাবুল আজান — বাবু ফজলিস-সালাতি ফি জামাআহ \\
\textbf{সনদের মান} & \textbf{সহীহ} (বুখারি-মুসলিম) \\
\hline
\end{tabular}
\end{table}

\#\#\#\# ২. পটভূমি (\lat{Context})

একাকী নামাজের তুলনায় জামাতের সওয়াব — ২৫ বা ২৭ গুণ; সংখ্যার পার্থক্য বর্ণনাকারীদের ভিন্নতার কারণে, মূল বার্তা এক।

\#\#\#\# ৩. পূর্ণ বর্ণনা (বাংলা, \lat{pure})

\begin{quote}
\textbf{\lat{“}জামাতে নামাজ একাকী নামাজের চেয়ে ২৭ গুণ (অন্য বর্ণনায় ২৫ গুণ) বেশি মর্যাদাপূর্ণ।\lat{”}}
\end{quote}

\#\#\#\# ৪. বিশ্লেষণ (\lat{Analysis})

\textbf{১.} মর্যাদার কারণ — জামাতে একতা, ইমামের অনুসরণ, কাতারের \textbf{\lat{discipline} (শৃঙ্খলা)}।
\textbf{২.} ২৫/২৭ — কোনো বর্ণনাকারীর স্মৃতিভেদ; উভয় বর্ণনাই সহীহ; আলেমরা এটিকে \lat{“}অন্তত ২৫, বেশি হলে ২৭\lat{”} বা তুলনার রূপক হিসেবে দেখেছেন।

\#\#\#\# ৫. শিক্ষা (\lat{Lessons})

\begin{enumerate}
\item পুরুষদের জন্য জামাত — সর্বোচ্চ চেষ্টা; মহিলারা ঘরে নামাজই উত্তম (আবু দাউদ ৫৭০)।
\item জামাতের সুযোগ হাতছাড়া না করা — ২৭ গুণের \textbf{\lat{opportunity} (সুযোগ)}।
\end{enumerate}
\medskip\hrule\medskip

\subsubsection*{হাদিস ১৮ — ঘরে নফল, ঘরকে কবর বানিও না (সহীহ বুখারি ৪৩২)}

\#\#\#\# ১. সূত্র কার্ড

\begin{table}[h]\centering\small
\begin{tabular}{ll}
বিষয় & বিবরণ \\
\hline
\textbf{বর্ণনাকারী} & আবদুল্লাহ ইবনে উমার (রাঃ) \\
\textbf{গ্রন্থ} & সহীহ বুখারি ৪৩২; সহীহ মুসলিম ৭৭৭ \\
\textbf{অধ্যায়} & কিতাবুস-সালাত — বাবু কিরাআতিল কুরআন ফিল বাইত \\
\textbf{সনদের মান} & \textbf{সহীহ} (বুখারি-মুসলিম) \\
\hline
\end{tabular}
\end{table}

\#\#\#\# ২. পটভূমি (\lat{Context})

নফল নামাজ ঘরে পড়ার নির্দেশ — যেন ঘর থেকে ইবাদতের আলো না হারায়।

\#\#\#\# ৩. পূর্ণ বর্ণনা (বাংলা, \lat{pure})

\begin{quote}
\textbf{\lat{“}তোমরা নিজেদের ঘরে (নফল) নামাজ পড়ো এবং ঘরগুলোকে কবরে পরিণত করো না।\lat{”}}
\end{quote}

\#\#\#\# ৪. বিশ্লেষণ (\lat{Analysis})

\textbf{১.} ফরজ নামাজ মসজিদে (জামাত), নফল নামাজ \textbf{ঘরে (\lat{at} \lat{home})} — উভয়ের সমন্বয়ই সুন্নত।
\textbf{২.} ঘরে নফল নামাজের বরকত — পরিবারের সামনে নামাজের \textbf{উদাহরণ (\lat{example})} স্থাপন।

\#\#\#\# ৫. শিক্ষা (\lat{Lessons})

\begin{enumerate}
\item সুন্নত/নফল নামাজ ঘরে পড়ার \textbf{\lat{habit} (অভ্যাস)} — অন্তত যুহরের আগের ৪ রাকাত।
\item ঘরের ইবাদত পরিবারের ঈমানের পরিবেশ তৈরি করে।
\end{enumerate}
\medskip\hrule\medskip

\subsection*{পার্ট ৩ — খুশু, মান ও সতর্কতা}

\subsubsection*{হাদিস ১৯ — ইতমিনানহীন নামাজ বাতিল (সহীহ বুখারি ৭৯১)}

\#\#\#\# ১. সূত্র কার্ড

\begin{table}[h]\centering\small
\begin{tabular}{ll}
বিষয় & বিবরণ \\
\hline
\textbf{বর্ণনাকারী} & হুযাইফা ইবনুল ইয়ামান (রাঃ) — জায়েদ ইবনে ওয়াহব থেকে \\
\textbf{গ্রন্থ} & সহীহ বুখারি ৭৯১ \\
\textbf{অধ্যায়} & কিতাবুল আজান — বাবু ইতমিনানিল মাসজুদ \\
\textbf{সনদের মান} & \textbf{সহীহ} (মাওকুফ — সাহাবির বক্তব্য, তবে এর বিধান নবীজির শিক্ষার সমর্থিত) \\
\hline
\end{tabular}
\end{table}

\#\#\#\# ২. পটভূমি (\lat{Context})

হুযাইফা (রাঃ) এক ব্যক্তিকে দেখলেন — রুকু-সিজদা ঠিকভাবে করছে না; তার প্রতিক্রিয়া নামাজের মানের গুরুত্ব বোঝায়।

\#\#\#\# ৩. পূর্ণ বর্ণনা (বাংলা, \lat{pure})

\begin{quote}
হুযাইফা (রাঃ) এমন এক ব্যক্তিকে দেখলেন, যে রুকু ও সিজদা ঠিকভাবে (ইতমিনানসহ) করছে না। তিনি বললেন: \\
\textbf{\lat{“}তুমি নামাজ পড়োনি! আর তুমি যদি মারা যাও, তবে তুমি মুহাম্মদ (সাল্লাল্লাহু আলাইহি ওয়া সাল্লাম)-এর সেই সুন্নতের উপর মারা গেলে না, যার উপর তিনি ছিলেন।\lat{”}}
\end{quote}

\#\#\#\# ৪. বিশ্লেষণ (\lat{Analysis})

\textbf{১.} ইতমিনান — রুকু, কাউমা, সিজদা, জালসায় শরীর স্থির রাখা; মুসি' সালাতুহু হাদিসে (বুখারি ৭৫৭, মুসলিম ৩৯৭) নবীজি এটিই শিখিয়েছেন।
\textbf{২.} দ্রুত ঠোকরানো নামাজ — নবীজির \textbf{সুন্নত (\lat{sunnah})} থেকে বিচ্যুতি; সবচেয়ে বড় চোর সেই, যে নিজের নামাজ চুরি করে (বুখারি ৭৯১-এর সংলগ্ন বর্ণনা)।

\#\#\#\# ৫. শিক্ষা (\lat{Lessons})

\begin{enumerate}
\item প্রতিটি রুকন — অন্তত একবার \lat{“}সুবহানা রাব্বিয়াল আজিম/আ'লা\lat{”} বলার মতো স্থির হোন।
\item দ্রুত নামাজের \textbf{\lat{habit} (অভ্যাস)} ভাঙুন — ধীরে, স্থিরভাবে।
\end{enumerate}
\medskip\hrule\medskip

\subsubsection*{হাদিস ২০ — নামাজের অল্পাংশই লেখা হয় (হাসান আবু দাউদ ৭৯৬)}

\#\#\#\# ১. সূত্র কার্ড

\begin{table}[h]\centering\small
\begin{tabular}{ll}
বিষয় & বিবরণ \\
\hline
\textbf{বর্ণনাকারী} & সাহাবি থেকে বর্ণিত (বর্ণনাটি বিভিন্ন সনদে এসেছে) \\
\textbf{গ্রন্থ} & সুনানে আবু দাউদ ৭৯৬ \\
\textbf{অধ্যায়} & কিতাবুস-সালাত — বাবু মা জাআ ফিমান লাম ইউতিম্মা রুকুআহু \\
\textbf{সনদের মান} & \textbf{হাসান} (আল-আলবানী) \\
\hline
\end{tabular}
\end{table}

\#\#\#\# ২. পটভূমি (\lat{Context})

নামাজ পড়েও মন নামাজে না থাকলে — নামাজের পুরো সওয়াব পাওয়া যায় না; নবীজি (সা.) এই বাস্তবতা তুলে ধরেছেন।

\#\#\#\# ৩. পূর্ণ বর্ণনা (বাংলা, \lat{pure})

\begin{quote}
\textbf{\lat{“}মানুষ নামাজ আদায় করে, কিন্তু তার নামাজ থেকে দশমাংশ, নবমাংশ, অষ্টমাংশ, সপ্তমাংশ, ষষ্ঠমাংশ, পঞ্চমাংশ — এমনকি তার অর্ধেকও লেখা হয় না।\lat{”}} (অর্থাৎ যতটুকু মনোযোগসহ পড়েছে, ততটুকুই মূল্য পায়)
\end{quote}

\#\#\#\# ৪. বিশ্লেষণ (\lat{Analysis})

\textbf{১.} সওয়াব নির্ভর করে \textbf{মনোযোগের (\lat{attention})} উপর — নামাজের পরিমাণ নয়।
\textbf{২.} খুশুর অভ্যাস গড়লে নামাজের পুরো সওয়াব — এটাই মুমিনের লক্ষ্য।

\#\#\#\# ৫. শিক্ষা (\lat{Lessons})

\begin{enumerate}
\item নামাজের সময় মোবাইল বন্ধ রাখুন — \textbf{\lat{distraction} (বিক্ষেপ)} সওয়াব কেড়ে নেয়।
\item অর্থ বুঝে নামাজ পড়ার চেষ্টা — মনোযোগের সবচেয়ে কার্যকর উপায়।
\end{enumerate}
\medskip\hrule\medskip

\subsubsection*{হাদিস ২১ — আযানের জবাব ও দোয়া (সহীহ বুখারি ৬১১, ৬১৪)}

\#\#\#\# ১. সূত্র কার্ড

\begin{table}[h]\centering\small
\begin{tabular}{ll}
বিষয় & বিবরণ \\
\hline
\textbf{বর্ণনাকারী} & আবু সাঈদ আল-খুদরী (রাঃ); জাবির ইবনে আবদিল্লাহ (রাঃ) \\
\textbf{গ্রন্থ} & সহীহ বুখারি ৬১১, ৬১৪; সহীহ মুসলিম ৩৮৩ \\
\textbf{অধ্যায়} & কিতাবুল আজান — বাবু মা ইউকালু 'ইন্দাল আজান \\
\textbf{সনদের মান} & \textbf{সহীহ} (বুখারি-মুসলিম) \\
\hline
\end{tabular}
\end{table}

\#\#\#\# ২. পটভূমি (\lat{Context})

আযান শুনে মুয়াজ্জিনের উত্তর দেওয়া — নামাজের আমন্ত্রণের প্রতি সাড়া; এর সাথে সুপারিশের দোয়া শেখানো হয়েছে।

\#\#\#\# ৩. পূর্ণ বর্ণনা (বাংলা, \lat{pure})

\begin{quote}
\textbf{\lat{“}যখন তোমরা আযান শুনো, তখন মুয়াজ্জিন যা বলে তোমরাও তাই বলো।\lat{”}} (বুখারি ৬১১, মুসলিম ৩৮৩) \\
আযানের পর নবীজি (সাল্লাল্লাহু আলাইহি ওয়া সাল্লাম) বলেছেন: \\
\textbf{\lat{“}যে ব্যক্তি আযানের পর এই দোয়া পড়বে — ‘আল্লাহুম্মা রাব্বা হাযিহিদ দাওয়াতিত তাম্মাহ, ওয়াস সালাতিল কায়িমাহ, আতি মুহাম্মাদানিল ওয়াসিলাতা ওয়াল ফাদিলাহ, ওয়াব'আসহু মাকামাম মাহমুদানিল্লাযী ওয়া'আত্তাহ’ — তার জন্য কিয়ামতের দিন আমার সুপারিশ ওয়াজিব হবে।\lat{”}} (বুখারি ৬১৪)
\end{quote}

\#\#\#\# ৪. বিশ্লেষণ (\lat{Analysis})

\textbf{১.} আযানের জবাব — ঈমানের \textbf{প্রতিক্রিয়া (\lat{response})}; নীরব থাকা সুন্নতের খিলাফ।
\textbf{২.} দোয়াটি মুখস্থ করা — অল্প আমলে বিশাল পুরস্কার (সুপারিশের প্রতিশ্রুতি)।

\#\#\#\# ৫. শিক্ষা (\lat{Lessons})

\begin{enumerate}
\item আযান শুনলে উত্তর দিন — কাজের মাঝেও।
\item আযানের দোয়াটি মুখস্থ করুন — \textbf{\lat{intercession} (সুপারিশ)}-এর সুযোগ।
\end{enumerate}
\medskip\hrule\medskip

\subsection*{সংক্ষিপ্ত সূত্র তালিকা (এই ফাইলের হাদিস)}

\begin{table}[h]\centering\small
\begin{tabular}{lll}
\# & বিষয় & উৎস \\
\hline
১ & প্রথম হিসাব নামাজ & তিরমিযী ৪১৩; আবু দাউদ ৮৬৪; নাসাঈ ৪৬৫ \\
২ & নামাজ ত্যাগ = কুফর সীমা & মুসলিম ৮২; তিরমিযী ২৬২০ \\
৩ & অঙ্গীকার: নামাজ & তিরমিযী ২৬২১; ইবনে মাজাহ ১০৭৯ \\
৪ & দ্বীনের স্তম্ভ & তিরমিযী ২৬১৬ \\
৫ & আসর ছুটলে পরিবার-ধন & বুখারি ৫৫২; মুসলিম ৬২৬ \\
৬ & মুনাফিকদের ভারী ইশা-ফজর & বুখারি ৬৫৭; মুসলিম ৬৫১ \\
৭ & পাঁচ নামাজ কাফফারা & মুসলিম ২৩৩ \\
৮ & চোখের শীতলতা & নাসাঈ ৩৯৩৯ \\
৯ & শেষ অসিয়ত & আবু দাউদ ৫১৫৬; ইবনে মাজাহ ২৬৯৮ \\
১০ & ফজরের ২ রাকাত & মুসলিম ৭২৫ \\
১১ & জামাতে ফজর-ইশা & মুসলিম ৬৫৬ \\
১২ & নামাজের চাবি & আবু দাউদ ৬১; তিরমিযী ৪; ইবনে মাজাহ ২৭৫ \\
১৩ & পবিত্রতা ছাড়া কবুল নয় & মুসলিম ২২৪; তিরমিযী ১ \\
১৪ & যেমন দেখেছ আমাকে & বুখারি ৬৩১ \\
১৫ & ফাতিহা ছাড়া নেই & বুখারি ৭৫৬; মুসলিম ৩৯৪ \\
১৬ & জমিন মসজিদ & বুখারি ৩৩৫; মুসলিম ৫২১ \\
১৭ & জামাত ২৫/২৭ গুণ & বুখারি ৬৪৫; মুসলিম ৬৫০ \\
১৮ & ঘরে নফল & বুখারি ৪৩২; মুসলিম ৭৭৭ \\
১৯ & ইতমিনান & বুখারি ৭৯১ \\
২০ & নামাজের অল্পাংশ & আবু দাউদ ৭৯৬ (হাসান) \\
২১ & আযানের জবাব ও দোয়া & বুখারি ৬১১, ৬১৪; মুসলিম ৩৮৩ \\
\hline
\end{tabular}
\end{table}


\end{document}
